\documentclass[a4paper,11pt,fleqn]{article}
\usepackage{Eval}


\newcommand{\chnum}{Chapitres 1 et 2}
\newcommand{\chtitle}{Méthodes d'analyse chimique}
\newcommand{\dtype}{Evaluation}

\begin{document}
\maketitle

\section*{Exercice 1 : Fraîcheur du lait}
\vspace{-1cm}
\begin{figure}[h!]
\addtolength{\belowcaptionskip}{2mm}
\caption*{\textbf{Document 1 : Qualité du lait frais et acide lactique}}
\begin{tabular}{m{13cm} m{6cm}}
	Le lait frais n'est que très légèrement acide, mais cette acidité peut se développer assez vite pour des raisons diverses. En particulier le lactose en quantité notable ($50~g \cdot L^{-1}$ environ) se transforme en acide lactique sous l'action des bactéries. Ainsi l'industrie laitière vérifie l'état de conservation d'un lait en mesurant son acidité totale en "équivalent d'acide lactique" exprimé en degré Dornic (°D). Le lait frais doit avoir, selon les normes en vigueur, une acidité inférieure à 18°D. & \includegraphics[scale=.1]{Ressources/acide lactique.png}\\
\end{tabular}
\textbf{Données :}
	\begin{itemize*}
	\item L'acide lactique a pour formule $C_3H_6O_3$ et le couple acide/base associé est : $CH_3CHOHCOOH$ / $CH_3CHOHCOO^{-}$
	\item (1°D correspond à 0,1 g d'acide lactique par litre de lait)
	\item $M(C_3H_6O_3) = 90,08~g \cdot mol^{-1}$
	\end{itemize*}
\end{figure}
\vspace{-1cm}
\begin{figure}[h!]
	\caption*{\textbf{Document 2 : Protocole de dosage du lait}}
	\begin{tabular}{m{11cm} m{6cm}}
	Pour réaliser le dosage du lait : 
	\begin{itemize*}
		\item Un volume de $20,0~mL$ de lait sont prélevés dans un bécher de $250~mL$.
		\item Ajouter environ $100,0~mL$ d'eau pour que le pH-mètre soit bien immergé
		\item Remplir une burette graduée avec une solution de soude ($Na^{+}_{(aq)} + HO^{-}_{(aq)}$) de concentration $C_B = 0,050~mol \cdot L^{-1}$.
		\item Verser lentement la solution d'hydroxyde de sodium sous agitation douce et relever les valeurs de pH
		\item Tracer la courbe $pH = f(V_B)$ puis déterminer le volume équivalent
	\end{itemize*} & \includegraphics[scale=.4]{Ressources/Figure_3.png}\\
	\end{tabular}
	\end{figure}
\begin{figure}[h!]
	\caption*{\textbf{Document 3 : Titrage de l'échantillon de lait}}
	\centering \includegraphics[scale=.7]{Ressources/Figure_2.png}
\end{figure}
\vspace{-1.5cm}
\subsection*{Questions :}
\begin{enumerate}
	\item Déterminer l'équation support du dosage
	\item Légendez le schéma du \textbf{Document 2}
	\item Définir l'équivalence
	\item Déterminer graphiquement l'équivalence en faisant figurer sur le \textbf{Document 3} la méthode employée
	\item Calculer la concentration d'acide lactique du lait testé
	\item L'échantillon est-il commercialisable ?
\end{enumerate}

\section*{Exercice 2 : Solution nettoyante pour lentilles}
Un produit utilisé pour le nettoyage des lentilles de contact contient, comme seule espèce ionique, du chlorure de sodium. Le fabricant indique : "chlorure de sodium : $0,85~g$ pour $100~mL$ de solution".\\
La conductivité $\sigma$ de solutions étalon de concentration en quantité de matière $C$ en chlorure de sodium est mesurée. Le graphique $\sigma = f(C)$ est donné ci-dessous :
\begin{center}
	\includegraphics[scale=.6]{Ressources/Figure_1.png}
\end{center}
La solution commerciale $S_0$ est diluée 10 fois. La conductivité de la solution diluée nommée $S$ est $\sigma_S = 1,8~mS \cdot cm^{-1}$.
\begin{enumerate}
	\item Déterminer la concentration $C_S$ en chlorure de sodium de la solution diluée $S$
	\item Calculer la concentration $C_0$ de la solution $S_0$
	\item En déduire sa concentration en masse $\gamma_0$
	\item À partir des indications de la notice, calculer la concentration en masse $\gamma_{notice}$ en chlorure de sodium de la solution commerciale
	\item Calculer l'écart relatif : $\frac{\lvert\gamma_{notice} - \gamma_0 \rvert}{\gamma_{notice}}$
	\item La concentration en masse $\gamma_0$ en chlorure de sodium de la solution $S_0$ satisfait elle au critère de qualité ?
\end{enumerate}
\textbf{Données :} \\
$M(NaCl)=58,5~g \cdot mol^{-1}$\\
Critère de qualité : le contrôle qualité est considéré comme satisfaisant si l'écart relatif est inférieur à 5\%

\section*{Exercice 3 : Détartrage}
Des détartrants à base d’acide lactique permettent de décomposer le tartre qui se déposent sur les résistances chauffantes.
La formule brute de l’acide lactique est $C_3H_6O_3$, et sa formule et sa masse molaire sont données dans l'\textbf{Exercice 1}.\\
Au cours de la préparation d’une solution détartrante, on dissout une masse $m = 54,1~g$ d’acide lactique pour $1,0~L$ de solution. Après agitation, le pH de la solution mesuré est 1,9.
\begin{enumerate}
\item Établir la formule semi-développée de l’acide lactique et entourer les groupes caractéristiques. Préciser à quelle famille appartient cette molécule.

\item Écrire l’équation de la réaction entre l’acide lactique et l’eau. (on pourra noter AH la formule de l’acide lactique)

\item En supposant que la réaction est totale, déterminer le pH attendu.

\item Justifier que la réaction n’est pas totale.
\end{enumerate}
Le "tartre" est un résidu solide constitué majoritairement de carbonate de calcium $CaCO_3$. La réaction avec un acide permet d'éliminer les dépôts.
\begin{enumerate}
\setcounter{enumi}{4}
\item En se basant sur les couples acide/base donnés, proposer une équation pour l'équation de détartrage.
\item Est il possible de suivre cette réaction par conductimétrie ?
\end{enumerate}
\textbf{Données : couples acide/base} \\
$C_3H_6O_3$ /$C_3H_5O_3^{-}$\\
$CO_2$ / $CO_3^{2-}$


\end{document}
